\chapter{Mocninné řady}%bez taylorovy rady, jen z diff rovnic
V této kapitole se naučíme sčítat nekonečně mnoho čísel a vyjadřovat funkce jako \uv{nekonečně dlouhé polynomy}. To nám umožní řešit diferenciální rovnice čistě symbolicky. Když nekonečný proces sčítání utneme po dostatečném počtu kroků, dostaneme aproximaci daného řešení, což nám umožní počítat funkce $\sin(x),\cos(x),\exp(x),\arctan(x), ...$ a tímpádem čísla $e,\pi, \ln(2),\sqrt{2},...$ Tuto aproximaci můžeme dosadit do složitějších diferenciálních rovnic a tím získat rovnice jednodušší, s tou nadějí, že zjednodušené rovnice se budou chovat podobně jako rovnice původní. Nakonci kapitoly vyřešíme legendární paradoxy z antického řecka, kterými si filosofové lámali hlavu 2000 let.
